\chapter{Boas práticas e padronização de código}
\label{cap:boas-praticas}

\begin{objetivos}
Ao final deste capítulo você deve ser capaz de:
\begin{itemize}
  \item aplicar uma convenção de nomes de variáveis e tags;
  \item estruturar o programa por área e por responsabilidade;
  \item situar os requisitos de segurança antes da lógica de funcionamento;
  \item reduzir o tempo de varredura sem sacrificar legibilidade;
  \item revisar um programa com um checklist objetivo.
\end{itemize}
\end{objetivos}

\section{Por que padronizar}

Um programa de CLP é lido muito mais vezes do que é escrito. Ele é lido na partida, no turno,
na manutenção corretiva e, anos depois, por quem não participou do projeto. Padronizar é reduzir o
tempo que cada uma dessas pessoas gasta para entender o que está escrito --- e é o que separa um
programa que funciona de um programa que \emph{continua} funcionando.

\index{boas práticas de programação}

\section{Convenção de nomes}

Regras mínimas, que valem mais do que qualquer padrão sofisticado:

\begin{enumerate}
  \item \textbf{nome pelo significado do processo}, não pelo endereço: \texttt{Bomba\_Agua\_Liga} e
        não \texttt{DO12};
  \item \textbf{prefixo por tipo de ponto}: evita confundir entrada, saída, bit interno e variável
        analógica;
  \item \textbf{sem acento e sem caractere especial}: o nome vai circular entre sistemas, planilhas e
        documentação;
  \item \textbf{verbo para comando, substantivo para estado}: \texttt{Liga}, \texttt{Desliga},
        \texttt{Em\_Operacao}, \texttt{Falha};
  \item \textbf{um nome por conceito}: se o mesmo sinal aparece com dois nomes no programa, alguém vai
        errar;
  \item \textbf{prefixo de área}: em planta com várias áreas, o nome começa pela área
        (\texttt{ENV\_} para envase, \texttt{PRE\_} para prensa).
\end{enumerate}

\begin{table}[H]
  \centering
  \caption{Prefixos sugeridos por tipo de ponto.}
  \label{tab:prefixos}
  \begin{tabular}{l l l}
    \toprule
    Prefixo & Tipo & Exemplo \\
    \midrule
    \texttt{DI\_} & entrada digital & \texttt{DI\_Porta\_Fechada} \\
    \texttt{AI\_} & entrada analógica & \texttt{AI\_Nivel\_Tanque} \\
    \texttt{DO\_} & saída digital & \texttt{DO\_Bomba\_Liga} \\
    \texttt{AO\_} & saída analógica & \texttt{AO\_Valvula\_Posicao} \\
    \texttt{B\_}  & bit interno / auxiliar & \texttt{B\_Ciclo\_Iniciado} \\
    \texttt{T\_}  & temporizador & \texttt{T\_Partida\_Bomba} \\
    \texttt{C\_}  & contador & \texttt{C\_Pecas\_Turno} \\
    \texttt{V\_}  & variável de cálculo & \texttt{V\_Vazao\_Filtrada} \\
    \bottomrule
  \end{tabular}
\end{table}

\section{Requisitos antes da lógica}

O desenvolvimento de qualquer sistema --- independentemente da técnica --- exige primeiro conhecer a
tecnologia, as possibilidades de ligação e as características dos elementos. E exige reconhecer que
equipamentos podem causar danos ao processo por defeito de peças, erro de programação ou erro de
instalação, podendo colocar vidas em risco.

Por isso, todo projeto deve ser conduzido de modo a atender, nesta ordem de prioridade:

\begin{enumerate}
  \item \textbf{segurança de pessoas};
  \item \textbf{segurança da instalação};
  \item \textbf{garantia de funcionamento};
  \item \textbf{facilidade de manutenção};
  \item \textbf{facilidade de operação};
  \item \textbf{custo}.
\end{enumerate}

Custo é o último critério, e essa ordem não é retórica: um projeto mais barato que falha no item 1 ou
2 não é um projeto viável.

\begin{figure}[H]
  \centering
  \includegraphics[width=\linewidth]{14-fluxo-desenvolvimento.pdf}
  \caption{Fluxo de elaboração do programa, dos requisitos à liberação para uso.}
  \label{fig:fluxo-desenvolvimento}
\end{figure}

Depois do ciclo básico esquematizado, entram as \emph{condições marginais} de funcionamento e
segurança: ciclo único ou contínuo; partida e parada; manual e automático; parada de emergência e
desbloqueio.

\section{Segurança dentro da lógica --- e fora dela}

A lógica programada deve implementar intertravamento, condição de falha segura e tratamento de falha.
Mas há um limite, e ele precisa estar claro para todo aluno:

\begin{table}[H]
  \centering
  \caption{O que pertence à lógica programada e o que não pertence.}
  \label{tab:limite-seguranca}
  \begin{tabular}{p{6.4cm} p{8.2cm}}
    \toprule
    Pode ser resolvido no programa & \textbf{Não} pode depender apenas do programa \\
    \midrule
    Intertravamento de processo: só abre a válvula se a pressão estiver abaixo do limite
      & Parada de emergência de pessoa: exige circuito de segurança dedicado ou controlador de
        segurança certificado \\
    Sequência de partida e de parada com verificação de condições
      & Detecção de presença em zona de risco: exige dispositivo de segurança (cortina de luz,
        relé de segurança) \\
    Diagnóstico, alarme e registro de falha
      & Proteção contra falha do próprio controlador (o programa não pode detectar sua própria
        ausência de execução) \\
    Retomada segura após falta de energia
      & Bloqueio de energização: exige meios externos ao software \\
    \bottomrule
  \end{tabular}
\end{table}

\begin{atencao}
Programa não é dispositivo de segurança. Uma lógica de emergência escrita em CLP comum, por mais
cuidadosa que seja, não cumpre função de segurança: se o controlador travar, a lógica de emergência
também trava. Use controlador de segurança ou circuito dedicado --- e trate isso na especificação
(Capítulo~\ref{cap:porte}), não na programação.
\end{atencao}

\section{Eficiência do ciclo de varredura}

O tempo de varredura depende do tamanho do programa e da quantidade de pontos
(Capítulo~\ref{cap:funcionamento}). Para reduzir sem perder clareza:

\begin{itemize}
  \item \textbf{use tarefas}: o que precisa ser rápido roda em tarefa periódica curta; diagnóstico e
        relatório ficam em tarefa lenta;
  \item \textbf{evite laços longos} dentro do ciclo: processamento pesado deve ir para tarefa
        própria, com período maior;
  \item \textbf{não duplique lógica}: cada regra existe uma vez e é chamada de onde for preciso;
  \item \textbf{cuidado com a comunicação} dentro do ciclo: troca de dados pela rede é justamente o
        que o \emph{watchdog} vigia;
  \item \textbf{meça, não estime}: comissione lendo o tempo de ciclo real, não o calculado.
\end{itemize}

\section{Checklist de revisão de código}

Use esta lista antes de liberar qualquer programa --- inclusive o do seu trabalho de conclusão.

\begin{table}[H]
  \centering
  \caption{Checklist de revisão.}
  \label{tab:checklist-codigo}
  \begin{tabular}{p{0.9cm} p{13.7cm}}
    \toprule
    \# & Verificação \\
    \midrule
    1 & Todos os nomes seguem a convenção e descrevem o processo \\
    2 & Cada POU tem uma responsabilidade única e um cabeçalho de documentação
        (Capítulo~\ref{cap:documentacao}) \\
    3 & Não há lógica duplicada em pontos diferentes do programa \\
    4 & Temporizadores e contadores têm instância única e base de tempo declarada \\
    5 & Variáveis que devem sobreviver à falta de energia estão declaradas como retentivas \\
    6 & Não há uso de escrita direta de E/S que contorne a memória imagem sem justificativa registrada \\
    7 & Condições de partida e de parada estão implementadas conforme o descritivo operacional \\
    8 & Intertravamentos e limites estão presentes e testados \\
    9 & O programa trata a primeira varredura e a retomada após falta de energia \\
    10 & Falha de sensor e de comunicação tem comportamento definido (e seguro) \\
    11 & O tempo de ciclo medido atende ao requisito do processo \\
    12 & Não há força (\emph{force}) ativa em saída \\
    13 & O programa foi testado conforme o roteiro do Capítulo~\ref{cap:simulacao} \\
    14 & A revisão final foi feita por alguém que não escreveu o código \\
    \bottomrule
  \end{tabular}
\end{table}

O item 14 é o mais ignorado e o mais barato. Quem escreveu não enxerga o próprio pressuposto errado;
outra pessoa enxerga em minutos.

\begin{pratica}
Pegue um programa de exercício seu e aplique o checklist da Tabela~\ref{tab:checklist-codigo}.
Depois:
(a) reescreva os nomes de todas as variáveis segundo a Tabela~\ref{tab:prefixos};
(b) identifique qualquer lógica duplicada e transforme em bloco de função reutilizável;
(c) registre quantos itens falharam e o que isso significa para quem herdar o programa.
\end{pratica}

\begin{exercicios}
\begin{enumerate}[label=\alph*)]
  \item Explique, com um exemplo, por que nomear variável pelo endereço físico cria problema na
        manutenção.
  \item Ordene, do mais para o menos prioritário, os seis requisitos de projeto e justifique as duas
        primeiras posições.
  \item Dê um exemplo de intertravamento que pode ficar no programa e um que \emph{não} pode.
        Explique a diferença.
  \item Um programa tem tempo de varredura alto. Cite três providências de projeto que reduzem o
        tempo sem cortar funcionalidade.
  \item Explique por que ``força'' ativa em saída é um risco e como evitá-lo.
  \item Escolha três itens do checklist e descreva o defeito concreto que cada um previne.
  \item Por que a revisão por outra pessoa encontra erros que o autor não encontra?
  \item Cite uma prática de padronização que reduz custo de migração de fabricante
        (Capítulo~\ref{cap:iec-versus-legado}) e uma que reduz custo de manutenção.
\end{enumerate}
\end{exercicios}
